\documentclass[12pt]{report}
\usepackage[margin=0.5in]{geometry}

\begin{document}


\emph{fridge} is a refrigerator design code written in the C language.  This manual 
describes the features in \emph{fridge}, the modes of operation available, and provides 
examples.


\chapter{Numerical Approach}

This chapter describes the approaches taken to generating the numerical solution in 
\emph{fridge}.  At a high level, the key components of the process involve 1) solving 
the flow network problem, 2) solving the transient temperature and water vapor problem, 
3) solving the discretized lumped parameter conduction problem.

\section{Flow Network Problem Setup}

The flow network problem is a description of a flow network with nodes and edges.  
Nodes are locations in the network which have no volume and perfectly mix the 
flow coming into them.  All components of the system are treated as edges - 
refrigerator volumes, ducts, heat exchangers, fans, and dampers.  Components are 
the only features of the system which have volumes and thus thermal mass.


\section{Ambient Heat Leak and Radiation}



\section{Moisture Management}


\section{Modeling Solid Temperature}

Two obects are allowed in \emph{fridge} to allow solid body temperatures to be 
captured:  Thermal Connections, and Thermal Masses.  For transient analysis, there 
is an important truly transient thermal inertia effect to capture for both of 
these.  To capture this, both of these are treated as 7-element discretized 
lumped parameter models as described here.

\subsection{Temperature Update Stability Condition}
Because \emph{fridge} is written with conserved variables in mind, the temperature update 
is done using Fourier's Law rather than the Heat Equation.  Both have stability 
conditions, and the one employed in \emph{fridge} is:

\begin{equation}
   \frac{dt\:k}{dx} < 625
\end{equation}

The code checks for this condition and can adjust the time stepping for each 
thermal connection such that this condition is met.  When this is necessary, 
a local dt is assigned to a given thermal connection which is a fraction of 
the global dt.  The result is that thermal connections may undergo 2, 3, 4, etc. 
local time steps for each global time step in order to stay in line with the 
global time.

\chapter{Defining the Configuration}





\chapter{Running in Steady Mode}


\section{Steady Mode Outputs}

\chapter{Running in Transient Mode}


\section{Transient Mode Outputs}

\chapter{Appendix A:  Examples}

Simple Fan Example

Full Refrigerator Example



\end{document}
